% ******************************* Thesis Appendix C ****************************

\chapter{Comprehensive Training Plan (Deliverable D4): USC-PIS}
\label{appendix:training_plan}

\section{Program Overview}

\textbf{Objectives:} The primary objective is to equip all users with the digital competency required to operate the fully digital USC-PIS, reducing patient record retrieval time to well under 1 second (capitalizing on the 127.81ms search benchmark).

\textbf{Delivery Methods:} Training will be executed via a hybrid delivery approach, using hands-on demonstration modules supplemented by the comprehensive User Manuals.

\section{System Adoption Strategy (Addressing Bottlenecks)}

Initial database audits revealed critical system underutilization (e.g., 3 medical records for 9 patients, 1 dental record, and 0 file uploads). To solve this adoption issue, training will incorporate:
\begin{itemize}
    \item \textbf{Workflow Integration:} Practical demonstrations on incorporating the software directly into the clinic's daily operational workflows.
    \item \textbf{Incentives \& Benefits:} Showcasing how digital record-keeping eliminates redundant paperwork and provides instant clinical insights.
\end{itemize}

\section{Module 1: Clinical Operations (For Doctors \& Nurses)}

\begin{itemize}
    \item \textbf{Patient Intake:} Staff will learn to utilize the \textbf{Application Retooling Filter Bar} to segment patient datasets in real-time by \textbf{Role, Program, Year Level, and Registration Period}. Instruction will explicitly cover the Academic Year and Semester logic (e.g., filtering 1st Semester from August 1 to December 31) to route targeted student groups.
    \item \textbf{Record Management \& Security:} Users will master the \textbf{Unified Record Component} for entering vitals and the mandatory "Concern / Reason for Visit" fields. A highly technical segment will focus on the \textbf{\texttt{pgcrypto} encryption architecture}, explaining how patient diagnoses and names are automatically hashed as secure \textbf{\texttt{BinaryField}} strings via backend signals.
\end{itemize}

\section{Module 2: Specialized Clinical Workflows (For Dentists)}

\begin{itemize}
    \item \textbf{Dental Charting \& Streamlined Entry:} Dentists will be instructed on charting tooth conditions using the \textbf{FDI notation (11-48)} numerical range, strictly enforced via regex. To accelerate high-volume clinic days, Dentists will master the \textbf{"Consultation Only" mode}, which makes full diagnoses optional, and utilize the explicit "Referral To" field.
\end{itemize}

\section{Module 3: Administrative Output (For Staff \& Admins)}

\begin{itemize}
    \item \textbf{Certificate Issuance:} Staff will learn the workflow for generating the \textbf{ACA-HSD-04F Medical Certificate} directly from a health record. Training will emphasize that while Nurses may draft the certificate, the system mandates secure approval executed solely by the \textbf{Doctor} before the PDF is rendered.
    \item \textbf{Reporting:} Staff will practice generating \textbf{Visit Trends} and \textbf{Treatment Outcome Reports} using the Celery-powered asynchronous reporting module, achieving data outputs in under 1 second without blocking the interface.
\end{itemize}

\section{Module 4: Student Engagement (For Students \& Faculty)}

\begin{itemize}
    \item \textbf{Personal Health Access \& Feedback:} Students will be trained on navigating their personalized \textbf{Student Medical Record (SMR)} via the \textbf{/health-insights} tab to view non-interactive timelines. They will also learn how to respond to automated post-visit emails by submitting evaluations and interacting with Health Campaigns.
\end{itemize}

\section{Assessment \& Evaluation}

\textbf{Training Success \& Adoption Metrics:} Rather than utilizing formal proficiency exams, the success of the training program and system adoption will be evaluated during the Pilot Deployment phase (targeting 70\% of 1st-year CpE students and selected clinic staff). The evaluation will be based on two realistic pillars:

\begin{itemize}
    \item \textbf{Direct User Feedback \& Issue Reporting:} Success will be measured by qualitative feedback gathered from the clinic staff and pilot students regarding the system's intuitiveness, as well as tracking any reported bugs or workflow bottlenecks.
    \item \textbf{Real-World System Performance:} Evaluating whether trained staff can execute operational tasks efficiently under real clinic loads, ensuring that benchmarks, such as the sub-500ms patient search latency and rapid PDF generation, are maintained during daily operations.
\end{itemize}
